% GENERATED FILE. Edit canonical lessons, then run npm run generate:books.
% canonical-source-hash: fnv1a64:c5ed3118fe53befb
% canonical-lessons: BN-C193-bishal, BN-C193-majhari, BN-C193-kachher, BN-C193-bktigoto, BN-C193-sorkari

\chapter{Huge, Medium, Nearby, Personal, Public}
\label{ch:bn-huge-medium-nearby-personal-public}
\hlchaptermodality{193}

\begin{chapteropening}
\textbf{By the end of this chapter:} \emph{I can say huge, medium, nearby, personal and public.}

Complete the last lesson of chapter 193: I can say huge, medium, nearby, personal and public.
\end{chapteropening}

\section[\texorpdfstring{\bn{বিশাল} (bishāl)}{বিশাল (bishāl)}]{\bn{বিশাল} (bishāl) — huge}

\label{lesson:BN-C193-bishal}



\begin{quote}
Before the new one: say the Bengali for separate, then the Bengali for the same.
\end{quote}

\subsection*{You'll want to know: \bn{বিশাল}}
\textbf{\bn{বিশাল}} — \emph{bishāl} — \textquotedblleft{}huge\textquotedblright{}.

\begin{cousinweb}
\textbf{\bn{বিশাল}} (\emph{bishāl}) is Sanskrit \emph{viśāla}, \textquotedblleft{}broad, great\textquotedblright{}, taken over whole as a learned loan, a tatsama.
\end{cousinweb}

\begin{grammarlens}[title={What you've built}]
One more word for this chapter.
\end{grammarlens}

\subsection*{Your turn}
\begin{itemize}
\raggedright
  \item \emph{Say it:} \emph{bishāl}
  \item \emph{Say it:} \emph{bishāl}, once more
  \item \emph{Say it:} say \emph{bishāl}
  \item \emph{Recall:} say the Bengali for separate, then the Bengali for the same, then say \emph{bishāl} again
\end{itemize}

\begin{tcolorbox}[breakable,colback=teal!4,colframe=teal!35!black,title={Before you move on}]
What does \bn{বিশাল} mean? (\textquotedblleft{}Huge\textquotedblright{}.) Say it once more.
\end{tcolorbox}
\section[\texorpdfstring{\bn{মাঝারি} (mājhāri)}{মাঝারি (mājhāri)}]{\bn{মাঝারি} (mājhāri) — medium}

\label{lesson:BN-C193-majhari}



\begin{quote}
Before the new one: say the Bengali for the same, then the Bengali for huge.
\end{quote}

\subsection*{You'll want to know: \bn{মাঝারি}}
\textbf{\bn{মাঝারি}} — \emph{mājhāri} — \textquotedblleft{}medium\textquotedblright{}.

\begin{cousinweb}
\textbf{\bn{মাঝারি}} (\emph{mājhāri}) is built on \emph{mājh}, \textquotedblleft{}middle\textquotedblright{}, from Sanskrit \emph{madhya}, \textquotedblleft{}middle\textquotedblright{} (CDIAL 9804).
\end{cousinweb}

\begin{grammarlens}[title={What you've built}]
One more word for this chapter.
\end{grammarlens}

\subsection*{Your turn}
\begin{itemize}
\raggedright
  \item \emph{Say it:} \emph{mājhāri}
  \item \emph{Say it:} \emph{mājhāri}, once more
  \item \emph{Say it:} say \emph{mājhāri}
  \item \emph{Recall:} say the Bengali for the same, then the Bengali for huge, then say \emph{mājhāri} again
\end{itemize}

\begin{tcolorbox}[breakable,colback=teal!4,colframe=teal!35!black,title={Before you move on}]
What does \bn{মাঝারি} mean? (\textquotedblleft{}Medium\textquotedblright{}.) Say it once more.
\end{tcolorbox}
\section[\texorpdfstring{\bn{কাছের} (kāchher)}{কাছের (kāchher)}]{\bn{কাছের} (kāchher) — nearby}

\label{lesson:BN-C193-kachher}



\begin{quote}
Before the new one: say the Bengali for huge, then the Bengali for medium.
\end{quote}

\subsection*{You'll want to know: \bn{কাছের}}
\textbf{\bn{কাছের}} — \emph{kāchher} — \textquotedblleft{}nearby\textquotedblright{}.

\begin{cousinweb}
\textbf{\bn{কাছের}} (\emph{kāchher}) is \emph{kāchh}, \textquotedblleft{}nearness\textquotedblright{}, + \emph{-er}, \textquotedblleft{}of\textquotedblright{}. CDIAL has no entry for \emph{kāchh}, but it appears to belong with Sanskrit \emph{kakṣa}, \textquotedblleft{}armpit, side\textquotedblright{}: there CDIAL lists Bengali \emph{kāṁkh}, \textquotedblleft{}flank\textquotedblright{}, Nepali \emph{kāch}, \textquotedblleft{}side, edge\textquotedblright{}, and Assamese \emph{kāxe}, \textquotedblleft{}by the side of\textquotedblright{}, and notes that some words for \textquotedblleft{}beside\textquotedblright{} probably grew from the sense \textquotedblleft{}side\textquotedblright{} (CDIAL 2588).
\end{cousinweb}

\begin{grammarlens}[title={What you've built}]
One more word for this chapter.
\end{grammarlens}

\subsection*{Your turn}
\begin{itemize}
\raggedright
  \item \emph{Say it:} \emph{kāchher}
  \item \emph{Say it:} \emph{kāchher}, once more
  \item \emph{Say it:} say \emph{kāchher}
  \item \emph{Recall:} say the Bengali for huge, then the Bengali for medium, then say \emph{kāchher} again
\end{itemize}

\begin{tcolorbox}[breakable,colback=teal!4,colframe=teal!35!black,title={Before you move on}]
What does \bn{কাছের} mean? (\textquotedblleft{}Nearby\textquotedblright{}.) Say it once more.
\end{tcolorbox}
\section[\texorpdfstring{\bn{ব্যক্তিগত} (bæktigôtô)}{ব্যক্তিগত (bæktigôtô)}]{\bn{ব্যক্তিগত} (bæktigôtô) — personal}

\label{lesson:BN-C193-bktigoto}



\begin{quote}
Before the new one: say the Bengali for medium, then the Bengali for nearby.
\end{quote}

\subsection*{You'll want to know: \bn{ব্যক্তিগত}}
\textbf{\bn{ব্যক্তিগত}} — \emph{bæktigôtô} — \textquotedblleft{}personal\textquotedblright{}.

\begin{cousinweb}
\textbf{\bn{ব্যক্তিগত}} (\emph{bæktigôtô}) is built from Sanskrit \emph{vyakti}, \textquotedblleft{}an individual\textquotedblright{}, and \emph{gata}, \textquotedblleft{}gone to, concerning\textquotedblright{}, both taken over whole: \textquotedblleft{}concerning the person\textquotedblright{}.
\end{cousinweb}

\begin{grammarlens}[title={What you've built}]
One more word for this chapter.
\end{grammarlens}

\subsection*{Your turn}
\begin{itemize}
\raggedright
  \item \emph{Say it:} \emph{bæktigôtô}
  \item \emph{Say it:} \emph{bæktigôtô}, once more
  \item \emph{Say it:} say \emph{bæktigôtô}
  \item \emph{Recall:} say the Bengali for medium, then the Bengali for nearby, then say \emph{bæktigôtô} again
\end{itemize}

\begin{tcolorbox}[breakable,colback=teal!4,colframe=teal!35!black,title={Before you move on}]
What does \bn{ব্যক্তিগত} mean? (\textquotedblleft{}Personal\textquotedblright{}.) Say it once more.
\end{tcolorbox}
\section[\texorpdfstring{\bn{সরকারি} (sôrkāri)}{সরকারি (sôrkāri)}]{\bn{সরকারি} (sôrkāri) — public, governmental}

\label{lesson:BN-C193-sorkari}



\begin{quote}
Before the new one: say the Bengali for nearby, then the Bengali for personal.
\end{quote}

\subsection*{You'll want to know: \bn{সরকারি}}
\textbf{\bn{সরকারি}} — \emph{sôrkāri} — \textquotedblleft{}public, governmental\textquotedblright{}.

\begin{cousinweb}
\textbf{\bn{সরকারি}} (\emph{sôrkāri}) is a loan, usually traced to Persian \emph{sarkār}, \textquotedblleft{}government; master\textquotedblright{}. It carries the adjective ending \emph{-i}.
\end{cousinweb}

\begin{grammarlens}[title={What you've built}]
One more word for this chapter.
\end{grammarlens}

\subsection*{Your turn}
\begin{itemize}
\raggedright
  \item \emph{Say it:} \emph{sôrkāri}
  \item \emph{Say it:} \emph{sôrkāri}, once more
  \item \emph{Say it:} say \emph{sôrkāri}
  \item \emph{Recall:} say the Bengali for nearby, then the Bengali for personal, then say \emph{sôrkāri} again
\end{itemize}

\begin{tcolorbox}[breakable,colback=teal!4,colframe=teal!35!black,title={Before you move on}]
What does \bn{সরকারি} mean? (\textquotedblleft{}Public, governmental\textquotedblright{}.) Say it once more.
\end{tcolorbox}
